\subsubsection{A Complete Two-Object Rollback Cascade}
\label{sec:tw-cascade}

The original three-car run does not force recovery. To see the original
machinery actually restore state, replay inputs, and cancel an output,
use the supplemental application in
\filepath{tutorial/examples/08-timewarp-rollback.abclr2}.
It introduces no new simulation algorithm: the unchanged original
loader, metaobjects, evaluator, environments, and history queues do the
recovery. The new application merely makes their effects easy to inspect.

The sink stores the latest received value and a reverse-order trace:

\wholelisting[firstline=6,lastline=12]{tutorial/examples/08-timewarp-rollback.abclr2}

The source adds each input to its total and sends that total to the sink
one simulated unit later:

\wholelisting[firstline=14,lastline=20]{tutorial/examples/08-timewarp-rollback.abclr2}

Neither application script contains an undo operation. Assignment goes
through the original environment protocol; the timed send is recorded
by the original evaluator. The trace uses \code{cons} to construct a new
list and rebinds \code{trace}, rather than destructively modifying a list
shared with earlier snapshots. That distinction matters for shallow
binding-history restoration.

\paragraph{Arrange a late event deliberately.}
Run \code{make -C tutorial timewarp-check}; the runner copies the
original directory and this supplemental file into
\filepath{build/timewarp-tests/}. For an interactive reproduction,
start a fresh interpreter in that prepared directory and load
\code{(inc "08-timewarp-rollback" :use-source t)}.
The example loads \code{tw-load} and resets the rollback monitor before
creating the two objects. It initially uses Lowest-LVT-first scheduling.

Inject a time-10 event, let the finite computation settle, and only
then inject a time-5 event:

\begin{codeblock}
\begin{lstlisting}
[[meta rollback-source] <= [:message [:add 10] nil nil 0 10]]
; Source total becomes 10 at LVT 10.
; Sink receives (:VALUE 10) at LVT 11.

[[meta rollback-source] <= [:message [:add 1] nil nil 0 5]]
; Source rolls back; its earlier output is cancelled at the sink.
; Final source total = 11, sink value = 11.
\end{lstlisting}
\end{codeblock}

These are deliberate meta-level event injections, not ordinary
application requests. The two \code{nil} fields are the initial reply
destination and sender; the final numbers are \code{VST = 0} and
\code{VRT = 10} or 5. No application reply is requested. The source's
subsequent timed sends name the source object as sender automatically.
The current interactive top-level drains runnable work after each
entered form. The checks exploit that boundary to ensure the first
output has \emph{already executed}, rather than merely waiting in the
sink's queue. This staged arrival order is a test arrangement, not a
promise that arbitrary concurrent clients see the same schedule.

\paragraph{Follow the histories, not just the final sum.}
After the first injection, the source has saved its initial frame at
time 0, set \code{LVT = 10}, and changed \code{total} from 0 to 10.
Its output history contains
\code{[rollback-sink [:value 10] nil 10 11]}.
The sink has likewise saved its initial frame at 0 before processing
time 11. Its value is 10 and its trace is \code{((11 10))}.

Now the late positive event has \code{VRT = 5 < LVT = 10}.
The source's \code{:begin-tw} handler takes the recovery branch.
It restores the saved frame at 0, so both its local time and total
return to zero. Its input queue has already dequeued the offending
time-5 event and the previous time-10 event. Rewinding to 5 makes both
available for replay, sorted by receive time. It must not discard the
late event after using it to trigger recovery.

The source extracts outputs whose send time is at least 5. The old
\code{[:value 10]} send qualifies because its send time was 10;
its receive time 11 is preserved in the cancellation envelope.
Here is the original anti-messenger's active handler:

\wholelisting[firstline=135,lastline=142]{examples/time-warp/tw-meta.abclr2}

The first pattern ends recursion on an empty history with \code{:fin}.
The second uses \code{bracket*} to separate one output record from the
remaining history. It sends directly to the target's metaobject,
changing the envelope's tag to \code{:anti-message}, but preserving
message, reply, sender identity, and both timestamps. Then it processes
the remaining records with the same \code{R}.

The source waits for the anti-messenger's completion before sending
itself \code{:end}. That completion means the cancellation requests
were issued, not that every receiving simulation object has completed
its own rollback. The target sends are asynchronous.
In this checked staged run, the sink receives:

\begin{codeblock}
\begin{lstlisting}
[:anti-message [:value 10] nil rollback-source 10 11]
\end{lstlisting}
\end{codeblock}

Because the positive event has already executed, it is in the sink's
\code{old} input history, not its pending queue. The anti-message
therefore cannot simply remove a pending positive event. The sink
dequeues the anti-message and rolls back to a frame strictly before
11: its initial frame at 0. This restores \code{value = 0} and
\code{trace = nil}. Rewinding the input history brings the old positive
and its anti-message together; \code{remove-anti} removes that pair,
so the invalid \code{[:value 10]} event is not replayed.
The sink has no timed outputs of its own to cancel in this application.

The source then processes its revised inputs at 5 and 10. At 5 it
adds 1 to zero and sends value 1 for receive time 6. At 10 it adds
10 to that corrected state and sends value 11 for receive time 11.
These are new outputs of replayed work, not the old output being
retained under another name. The sink's revised history is therefore:

\begin{center}
\begin{tabular}{rll}
\textbf{Receive time} & \textbf{Application event} & \textbf{Resulting sink value}\\\hline
6 & \code{[:value 1]} & 1\\
11 & \code{[:value 11]} & 11\\
\end{tabular}
\end{center}

The final reverse trace is \code{((11 11) (6 1))}; the speculative
\code{(11 10)} entry has disappeared. The source's final local time is
10, the sink's is 11, and the shared monitor records two rollbacks:
one caused by a late positive input, one caused by cancellation of
already processed output. Console cancellation and rollback lines
remain visible: restoring simulation bindings does not erase printing.

\paragraph{Inspect without injecting another application event.}
The original \code{TW-meta} does not implement an application-state
query. Its own ordinary metaobject can expose the simulation machinery's
bindings. For example:

\begin{codeblock}
\begin{lstlisting}
[source-meta-state := [[meta [meta rollback-source]] <== :state]]
[source-state := [source-meta-state <== [:value-of 'state]]]
[source-state <== [:value-of 'total]] ; => 11
[source-meta-state <== [:value-of 'LVT]] ; => 10
\end{lstlisting}
\end{codeblock}

The inner \code{meta} selects the simulation object's interpreter;
the outer selects that interpreter's ordinary metaobject.
Its \code{:state} reply is an environment whose binding named
\code{state} points to the undoable \emph{application} environment.
\code{LVT}, \code{pqueue}, and \code{output-history} instead belong
to the simulation metaobject's own frame. Do not confuse these two
state environments or treat the returned objects as copied snapshots.

\paragraph{Evidence and limits.}
Twelve assertions run for this cascade with each evaluator scheduler,
in separate fresh processes: Lowest-LVT-first and FIFO. They check
the speculative values and output record, corrected values and trace,
both final local times, the replacement output history, the replayed
source input history, the sink's retained corrected inputs, empty
pending event queues, and the two-rollback counter. Transcript
assertions require the cancellation and both restoration reports.
The finite result is the same under these two checked arrangements;
the test is not a performance comparison.

This closes the gap between isolated history tests and a concrete
metaobject/evaluator recovery cascade. It does not certify every
cancellation arrival order, simultaneous or duplicate event,
zero-time rollback boundary, or timestamped now-type send. In
particular, output-history matching is based on the original event
fields rather than unique event identifiers. The same original
\code{remove-anti} helper compares envelope tails, not the tags;
do not infer a general duplicate-delivery policy from this one pair.
The experimental NR/NS evaluator variants remain separate cases.

\paragraph{Check your understanding.}
Why is restoring only the source's total insufficient? Its speculative
output has already changed the sink. Why is the final value 11 rather
than 1? The time-10 input is replayed after the new time-5 input.
Why does the sink not end with both time-11 values in its trace?
The old input is cancelled during rewind and its earlier state effect
is rolled back. Does waiting for the anti-messenger join the entire
cascade? No; it joins cancellation issuance, not all recipient work.
